\section{State-of-the-art}
\label{sec:state-of-the-art}
The intersection of large language models, retrieval-augmented generation, and personalised education has seen rapid advancement in recent years. This section surveys the most recent and relevant works that define the current state of the art.


\subsection{GraphRAG: From Local to Global Retrieval}

Edge et al.~\cite{edge2024localglobalgraphrag} propose GraphRAG, a graph-based retrieval-augmented generation framework that addresses the limitations of conventional RAG systems when answering queries requiring information distributed across a large document collection. Traditional RAG architectures retrieve a limited set of text chunks according to their semantic similarity to the input query. Although this retrieval strategy is effective for fact-oriented questions, it frequently fails to capture high-level themes and long-range relationships that span multiple documents. Consequently, responses generated by standard RAG systems may be incomplete when the required information is dispersed throughout the corpus.

To overcome this limitation, GraphRAG constructs a structured knowledge graph during the indexing phase. Rather than representing the corpus as an independent collection of text chunks, the framework employs a large language model to extract entities and semantic relationships from the documents. These entities are represented as graph nodes, while their relationships form the graph edges. Subsequently, a community detection algorithm partitions the graph into densely connected subgraphs representing coherent semantic topics. For each community, an LLM generates a concise summary describing the principal concepts and relationships contained within that community. The resulting hierarchical representation provides both fine-grained entity-level information and higher-level semantic summaries.

During inference, GraphRAG supports two complementary retrieval strategies. For local queries concerning specific entities or concepts, the system retrieves the relevant subgraph together with its neighbouring nodes. Conversely, for global queries that require an understanding of the document collection as a whole, GraphRAG retrieves multiple community summaries and synthesises them into a comprehensive response. This hierarchical retrieval mechanism enables the language model to reason over semantically organised information rather than isolated document fragments, thereby improving the completeness and coherence of generated answers.

Compared with conventional retrieval-augmented generation, GraphRAG offers several important advantages. The explicit graph representation captures semantic relationships that are not preserved by vector similarity alone, while the hierarchical community summaries facilitate corpus-level reasoning without requiring exhaustive retrieval of individual documents. As a result, the framework demonstrates improved performance on complex knowledge-intensive tasks, particularly those involving query-focused summarisation and multi-document reasoning.

The GraphRAG framework is closely related to the objectives of the present thesis. Both approaches employ structured knowledge representations to enhance retrieval beyond conventional embedding-based similarity search. However, whereas GraphRAG focuses primarily on improving information retrieval and summarisation across large document collections, the approach proposed in this thesis extends this concept to personalised educational question answering. Specifically, the proposed Tripartite RAFT framework integrates learner-specific information obtained from the learner model with structured knowledge retrieved from MathHub. Consequently, retrieval and answer generation are conditioned not only on the semantic relevance of the retrieved content but also on the student's current level of understanding and prerequisite knowledge, enabling explanations that are adapted to individual learners.


\subsection{Tripartite-GraphRAG via Plugin Ontologies}

This work \cite{banf2025tripartitegraphragpluginontologies} introduces Tripartite-GraphRAG for medical domain, a retrieval-augmented generation framework that addresses two persistent limitations in standard RAG pipelines: the reliance on embedding-based similarity for chunk selection, which produces noisy or incomplete prompts, and the scalability challenges of entity-based knowledge graph construction, such as entity resolution and deduplication.

The core contribution is a three-node knowledge graph connecting domain-specific objects of investigation (e.g., a patient anamnesis), a curated plugin ontology of concepts (e.g., clinical terms), and text chunks from source documents. At construction time, an LLM performs a concept-anchored pre-analysis of each text chunk, extracting and storing only information relevant to each concept as edge properties — an "information-preserving pre-compression" that discards textual noise while retaining granular, query-relevant content. This design avoids entity resolution entirely because the ontology provides stable, predefined concept nodes rather than extracted entities.

For inference, the authors formulate prompt assembly as an unsupervised node classification problem on a transformed graph. Each candidate text chunk is scored via cosine similarity between the object–concept and concept–chunk edge embeddings, and inclusion is decided against a concept-specific empirical distribution $P_c$ of these scores, controlled by hyperparameter $\alpha$. A secondary parameter $\beta$ allows co-occurring concepts within the same text chunk to influence each other's selection, inspired by Markov Random Field principles.

This method inspires the present work in the following way: we conceptualise the personalised information of each student as a plug-and-play component within the prompt system, while the remaining parts of the content — namely the MathHub repository and the trunking system — remain identical for every student. This design principle played a significant role in shaping this thesis.



\subsection{Guided Tours in ALeA}

Betzendahl et al.~\cite{10.1007/978-3-031-50485-3_39} introduce \emph{Guided Tours}, a core component of the Adaptive Learning Assistant (ALeA), an intelligent tutoring system designed to provide personalized educational dialogues based on an individual learner's knowledge and learning progress. Unlike conventional e-learning systems that present identical learning materials to all users, ALeA dynamically assembles educational interactions according to the learner's current competencies, prior experience, and learning objectives. The primary objective is to emulate the adaptive behaviour of a human tutor by presenting explanations and learning resources that are appropriate for each learner's individual background.

The personalization capabilities of ALeA are founded on four complementary models: a \emph{domain model}, a \emph{learner model}, a \emph{formulation model}, and a \emph{didactic model}. The domain model represents the concepts within the subject area and the semantic relationships between them, while the learner model maintains an estimate of the learner's competence for each concept. The formulation model stores alternative representations of educational content, including textual explanations, examples, exercises, and multimedia resources. Finally, the didactic model describes the pedagogical role of each learning object, such as whether it serves as a definition, example, theorem, exercise, or assessment item. Together, these models enable the system to select educational content that is both pedagogically appropriate and aligned with the learner's current level of understanding.

A distinguishing characteristic of ALeA is its reliance on semantically annotated learning materials. Course documents are enriched with structural and semantic metadata using the \texttt{sTeX} framework, allowing mathematical concepts, prerequisite relationships, learning objectives, and rhetorical classifications to be extracted automatically. The resulting semantic representation forms a knowledge graph that serves as the foundation for educational services such as guided tours, adaptive navigation, and learner modelling.

A guided tour is generated dynamically rather than being predefined. When a learner requests assistance on a particular concept, the system first analyses the learner model to estimate the learner's current level of understanding. Based on this assessment, prerequisite concepts are identified, appropriate learning objects are selected from the semantic repository, and an educational dialogue is assembled that gradually introduces the required knowledge before addressing the target concept. Throughout the interaction, learner responses are continuously used to update the learner model, allowing subsequent recommendations to reflect the learner's evolving knowledge state.

This work provides a solid foundation demonstrating how personalisation can influence responses within the educational domain. However, the system is entirely conversation-driven: if the user does not engage in dialogue or respond to its prompts, the system cannot be of use. In contrast, the goal of this thesis is to enable a student to pose questions and receive answers that are personalised to their individual needs, thereby addressing the "one-size-fits-all" problem.